% Generated draft: requires mathematical and source review.
\documentclass[11pt]{article}
\usepackage[margin=25mm]{geometry}
\usepackage{fontspec}
\setmainfont{Latin Modern Roman}
\usepackage{amsmath,amssymb,mathtools,unicode-math}
\setmathfont{Latin Modern Math}
\usepackage{enumitem,needspace,xurl,hyperref}
\setlength{\parindent}{0pt}
\setlength{\parskip}{5pt}
\setlength{\emergencystretch}{3em}
\newcommand{\N}{\mathbb{N}}
\newcommand{\Z}{\mathbb{Z}}
\newcommand{\Q}{\mathbb{Q}}
\newcommand{\R}{\mathbb{R}}
\newcommand{\C}{\mathbb{C}}
\newcommand{\F}{\mathbb{F}}
\newcommand{\A}{\mathbb{A}}
\newcommand{\PP}{\mathbb{P}}
\newcommand{\E}{\mathbb{E}}
\newcommand{\eps}{\varepsilon}
\newcommand{\dd}{\,\mathrm{d}}
\begin{document}
\section*{UnsolvedMath ID: 20000334}
\textbf{Source problem code:} catalogue entry
\subsection{Definitions and mathematical statement}
Morrison's AIM problem asks for easier estimates of Hilbert-point weights using the combinatorics of monomial subspaces of $H^0(X,L^m)$ for a polarized variety $(X,L)$ and a one-parameter subgroup $\lambda$. It asks in particular for more flexible filtrations with geometrically useful stages and their leading weight estimates, and for estimates built from nonnested collections of subspaces of $H^0(X,L)$. The exact Hilbert--Mumford sign convention, embedding, and notation $r_{\lambda,X}$ are not defined in the excerpt. The example about cusps is contextual and cannot be treated as an embedding-independent instability theorem.

To state the supplied report's linear-algebra result, let $V$ be a finite-dimensional filtered vector space with weight $w(v)$ the least filtration index containing a nonzero vector $v$. Suppose subspaces $U_1,\ldots,U_s$ have certified caps $w(v)\leq c_i$ for all nonzero $v\in U_i$, ordered $c_1\leq\cdots\leq c_s$, and a fallback cap $c_*\geq c_s$ holds on $V$. Put
\[
 R_0=0,\qquad R_i=\dim(U_1+\cdots+U_i),\qquad N=\dim V.
\]
The report's sharp cumulative-rank bound on minimum basis weight is
\[
 \min_{\mathcal B\text{ a basis of }V}\sum_{v\in\mathcal B}w(v)
 \leq\sum_{i=1}^s c_i(R_i-R_{i-1})+c_*(N-R_s).
\]
Sharp means optimal information from these caps and cumulative ranks alone, not equality for every geometric filtration. To convert this to $r_{\lambda,X}$ one must verify the source's weight/sign convention.

For $V_m=H^0(X,L^m)$ on an $n$-fold, if $R_i(m)=a_i m^n+O(m^{n-1})$, $N(m)=a m^n+O(m^{n-1})$, and caps $c_i(m)=b_i m+O(1)$, $c_*(m)=b_*m+O(1)$ have an eventual common ordering, the bound has leading coefficient
\[
 \sum_i b_i(a_i-a_{i-1})+b_*(a-a_s)
\]
in degree $m^{n+1}$, with $a_0=0$.

The report says a jet-surjectivity argument on smooth curves supplies the geometric Span Lemma used by Swinarski, replacing an unrestricted inclusion--exclusion identity that fails for arbitrary triples of subspaces. One explicit sufficient jet setting is surjectivity of $H^0(C,M)\to H^0(C,M|_{D_{\max}})$ for effective divisors $D_i$ and coefficientwise maximum $D_{\max}$; the subspaces $H^0(C,M(-D_i))$ then have their spans controlled by the corresponding coordinatewise jet conditions. The exact Span Lemma and degrees used in the report are not supplied.

Finally its two-point divisor family reportedly improves the leading estimate by $(b-a)(\beta-\gamma)m^2+O(m)$, conditional on a common low-weight cap and a fallback cap. The parameters and cap hypotheses are not defined in the record; strict improvement requires the indicated product positive. This conditional example is not a universal gain for arbitrary subspace families.
\subsection{Short English statement}
Improve Hilbert-weight estimates using geometric filtrations and nonnested subspaces. The report bounds minimum basis weight by ordered caps and cumulative span ranks, then extracts leading coefficients. It also replaces an invalid general inclusion--exclusion step by a jet argument on curves and gives a conditional two-point improvement whose detailed hypotheses are missing.
\subsection{Sources}
\begin{itemize}
\item Local registry: \url{data/unsolvedmath/problems.json}, permanent ID $20000334$, source code \texttt{catalogue entry}; Morrison's original problem and complete research summary.
\item AIM workshop ``Compact moduli spaces and birational geometry'', ``Open Problems in Compact Moduli Spaces and Birational Geometry'': \url{https://aimath.org/WWN/birational/birational.pdf}. Canonical location \texttt{aim-algebraic-geometry-notes.json}, \texttt{notes[333]}, source item 3 with extraction-broken numbering. The supplied verification excludes the appended heading ``4. Examples and applications''. No PDF was opened here.
\item The report names Swinarski's method and Span Lemma without a full bibliographic citation or exact statement; original remarks and literature field empty.
\item Supplied machine-generated research attempt 1: partial result; problem partially solved at the run and in the registry; rigor metadata ``proved'', novelty confidence low. The non-peer-reviewed bounds and corrections require independent verification before citation.
\end{itemize}
\end{document}
